\practicechapter{第 15 章实践：可信性能测量基础}

本章对应第一册第 15 章“可信性能测量基础”。正文讲时间源、warmup、重复测量、噪声、偏差、PMU 和报告结构。本章把这些原则落到当前项目的 Google Benchmark smoke：它只测解析 synthetic ELF 的路径，不测文件读取，不测打印，也不等于真实 workload 性能结论。

\practicesection{一、对应原书问题：benchmark 不是时间数字}

一个 benchmark 输出了纳秒数，并不自动可信。先问四个问题：测的输入是什么，计时边界在哪里，正确性是否已由测试保证，当前二进制身份是否固定。本项目里，正确性主要由 GTest/GMock 固定，benchmark 只对解析路径做 smoke 检查。

这就是“测量基础”的核心：性能数字必须附着在实验协议上。没有协议，数字无法复查；没有正确性测试，速度没有意义；没有 checksum 和构建信息，数字无法绑定到具体产物。

\practicesection{二、从特例推到规律：计时边界会污染结论}

特例一：如果 benchmark 循环里每次都构造 synthetic ELF，那么测到的是构造和分配成本，不只是解析成本。特例二：如果循环里打印 report，那么 I/O 会主导时间。特例三：如果 benchmark 读真实文件，磁盘缓存、文件系统和路径解析会进入计时边界。

由此推出规律：microbenchmark 要把 setup 放在计时边界外，把热路径放在计时边界内，并把这个边界写进报告。本项目 benchmark 选择 synthetic ELF，是为了让输入稳定、体积小、可重复。它不能代表所有真实 ELF 文件，但可以作为解析入口没有明显退化的 smoke。

\practicesection{三、实践任务：运行 benchmark 并解释限制}

先运行完整测试：

\begin{lstlisting}[language=bash]
cmake -S books/cpu-volume-1-practice \
      -B books/cpu-volume-1-practice/build/executable-evidence-debug \
      -DCMAKE_BUILD_TYPE=Debug
cmake --build books/cpu-volume-1-practice/build/executable-evidence-debug
ctest --test-dir books/cpu-volume-1-practice/build/executable-evidence-debug \
      --output-on-failure
\end{lstlisting}

然后找到 benchmark 可执行文件并运行。若本地 Google Benchmark 输出格式不同，以实际构建目录为准：

\begin{lstlisting}[language=bash]
books/cpu-volume-1-practice/build/executable-evidence-debug/labs/executable_evidence/cpu1ee_bench
\end{lstlisting}

阅读 \filepath{bench/executable\_evidence\_bench.cpp}，说明输入构造在哪里，计时循环在哪里，\code{benchmark::DoNotOptimize} 保护的是什么。报告中不要只贴结果数字，必须写清“这个 benchmark 没有测文件读取、没有测 CLI 参数解析、没有测打印、没有测真实大文件”。

\practicesection{四、验收：性能结论必须可降级}

本章交付物是 \filepath{reports/ch15-measurement-foundations.md}。通过标准如下：

\begin{enumerate}
  \item 报告记录构建类型、CPU 环境能获得的信息、benchmark 命令。
  \item 写清 benchmark 的输入和计时边界。
  \item 写清 GTest/GMock 与 benchmark 的分工：测试证明行为，benchmark 观察成本。
  \item 至少列出三个没有被当前 benchmark 覆盖的成本。
  \item 若结果波动或环境受限，结论要降级为 smoke，不写成确定性能排名。
\end{enumerate}

这章完成后，读者应该不再把 benchmark 当作“跑一下看数字”，而是当作一份需要协议、边界和限制的工程证据。
